\let\negmedspace\undefined
\let\negthickspace\undefined
\documentclass[journal]{IEEEtran}
\usepackage[a5paper, margin=10mm, onecolumn]{geometry}
%\usepackage{lmodern} % Ensure lmodern is loaded for pdflatex
\usepackage{tfrupee} % Include tfrupee package

\setlength{\headheight}{1cm} % Set the height of the header box
\setlength{\headsep}{0mm}     % Set the distance between the header box and the top of the text

\usepackage{gvv-book}
\usepackage{gvv}
\usepackage{cite}
\usepackage{amsmath,amssymb,amsfonts,amsthm}
\usepackage{algorithmic}
\usepackage{graphicx}
\usepackage{textcomp}
\usepackage{xcolor}
%\usepackage{txfonts}
\usepackage{listings}
\usepackage{enumitem}
\usepackage{mathtools}
\usepackage{gensymb}
\usepackage{comment}
\usepackage[breaklinks=true]{hyperref}
\usepackage{tkz-euclide} 
\usepackage{listings}
% \usepackage{gvv}                                        
\def\inputGnumericTable{}                                 
\usepackage[latin1]{inputenc}                                
\usepackage{color}                                            
\usepackage{array}                                            
\usepackage{longtable}                                       
\usepackage{calc}                                             
\usepackage{multirow}                                         
\usepackage{hhline}                                           
\usepackage{ifthen}                                           
\usepackage{lscape}
\usepackage{circuitikz}
\tikzstyle{block} = [rectangle, draw, fill=blue!20, 
    text width=4em, text centered, rounded corners, minimum height=3em]
\tikzstyle{sum} = [draw, fill=blue!10, circle, minimum size=1cm, node distance=1.5cm]
\tikzstyle{input} = [coordinate]
\tikzstyle{output} = [coordinate]


\begin{document}

\bibliographystyle{IEEEtran}
\vspace{3cm}

\title{4.13.99}
\author{AI25BTECH11030 -Sarvesh Tamgade}
{\let\newpage\relax\maketitle}

%\renewcommand{\thefigure}{\theenumi}
%\renewcommand{\thetable}{\theenumi}
\setlength{\intextsep}{10pt} 


\numberwithin{equation}{enumi}
\numberwithin{figure}{enumi}
%\renewcommand{\thetable}{\theenumi}


\textbf{Question}: Two planes are given:
\begin{align}
P_{1}:\;10x + 15y + 12z - 60 &= 0,\\
P_{2}:-2x + 5y + 4z - 20 &= 0.
\end{align}
Which of the following straight lines can be an edge of some tetrahedron whose two faces lie on $P_{1}$ and $P_{2}$?
\begin{enumerate}[label=(\alph*)]
  \item $\displaystyle \frac{x-1}{0} = \frac{y-1}{-5} = \frac{z-1}{4}$
  \item $\displaystyle \frac{x-6}{-5} = \frac{y-2}{3} = \frac{z}{-5}$
  \item $\displaystyle \frac{x}{-2} = \frac{y-4}{5} = \frac{z}{4}$
  \item $\displaystyle \frac{x}{1} = \frac{y-4}{-2} = \frac{z}{3}$
\end{enumerate}

\textbf{Solution:}\\
The two planes are:
\begin{align}
P_{1}:&\;\vec{n}_{1}^\top \vec{r} = d_{1}, 
\quad \vec{n}_{1} = \myvec{10 \\ 15 \\ 12},\; d_{1} = 60,\\
P_{2}:&\;\vec{n}_{2}^\top \vec{r} = d_{2}, 
\quad \vec{n}_{2} = \myvec{-2 \\ 5 \\ 4},\; d_{2} = 20.
\end{align}
A general line is
\begin{align}
\vec{r}(t) = \vec{r}_{0} + t\,\vec{d}.
\end{align}
For a line to be an edge of a tetrahedron whose two faces lie on $P_1$ and $P_2$, it must satisfy one of the conditions:
\begin{align}
t_1 &= \frac{d_1 - \vec{n}_1^\top\vec{r}_0}{\vec{n}_1^\top\vec{d}},&
t_2 &= \frac{d_2 - \vec{n}_2^\top\vec{r}_0}{\vec{n}_2^\top\vec{d}},
\end{align}
and then either
\begin{enumerate}
  \item $\vec{n}_1^\top\vec{r}_0=d_1$ and $\vec{n}_1^\top\vec{d}=0$ i.e. the line lies on plane $P_{1}$,or
  \item $\vec{n}_2^\top\vec{r}_0=d_2$ and $\vec{n}_2^\top\vec{d}=0$ i.e. the line lies on plane $P_{2}$,or
  \item $t_1\neq t_2$ i.e. the line intersects both planes on two different points.
\end{enumerate}
We test the options:
\begin{enumerate}[label=(\alph*)]
  \item $\displaystyle \vec{r}_0=\myvec{1\\1\\1},\,\vec{d}=\myvec{0\\-5\\4}$. 
    \(\vec{n}_1^\top\vec{r}_0=37,\;\vec{n}_2^\top\vec{r}_0=7,\;
    \vec{n}_1^\top\vec{d}=-27,\;\vec{n}_2^\top\vec{d}=-9,\)
    \[
      t_1=\tfrac{60-37}{-27}=-\tfrac{23}{27},\;
      t_2=\tfrac{20-7}{-9}=-\tfrac{13}{9},\quad t_1\neq t_2\ (\text{valid}).
    \]
  \item $\displaystyle \vec{r}_0=\myvec{6\\2\\0},\,\vec{d}=\myvec{-5\\3\\-5}$. 
    \(\vec{n}_1^\top\vec{r}_0=90,\;\vec{n}_2^\top\vec{r}_0=-2,\;
    \vec{n}_1^\top\vec{d}=-65,\;\vec{n}_2^\top\vec{d}=5,\)
    \[
      t_1=\tfrac{60-90}{-65}=\tfrac{6}{13},\;
      t_2=\tfrac{20-(-2)}{5}=\tfrac{22}{5},\quad t_1\neq t_2\ (\text{valid}).
    \]
  \item $\displaystyle \vec{r}_0=\myvec{0\\4\\0},\,\vec{d}=\myvec{-2\\5\\4}$. 
    \(\vec{n}_1^\top\vec{r}_0=60,\;\vec{n}_2^\top\vec{r}_0=20,\;
    \vec{n}_1^\top\vec{d}=103,\;\vec{n}_2^\top\vec{d}=45,\)
    \[
      t_1=\tfrac{60-60}{103}=0,\;
      t_2=\tfrac{20-20}{45}=0,\quad t_1=t_2\ (\text{invalid}).
    \]
  \item $\displaystyle \vec{r}_0=\myvec{0\\4\\0},\,\vec{d}=\myvec{1\\-2\\3}$. 
    \(\vec{n}_2^\top\vec{r}_0=20,\;\vec{n}_2^\top\vec{d}=0,\)
    \[
      \myvec{-2 \ 5 \ 4 }\vec{r}=20\quad\forall\,t,
      \quad\text{(the line lies on }P_2\text{) valid}.
    \]
\end{enumerate}
Therefore, the valid edges are \textbf{(A), (B), and (D)}.

\begin{figure}[htbp]
    \centering
    \includegraphics[width=0.65\linewidth]{FIG/fig1.png}
    \caption{Vector Representation}
    \label{fig:FIG/fig1.png}
    \end{figure}

\end{document}  
